\chapter{\texorpdfstring{$\beta$}{Beta}-eliminación}\label{ch:b1:elimination}

Trátese el 2-bromo-2-metilpropano con etóxido de sodio en etanol a
temperatura ambiente y el producto principal es un éter; caliéntese la
misma mezcla y el producto principal es un alqueno, el 2-metilpropeno, sin
ningún grupo etoxi. El ion etóxido no ha actuado como \omterm{def:b1:electronic-effects:nucleophile}{nucleófilo}, atacando
al carbono, sino como base, arrancando un protón del carbono vecino
mientras se iba el bromuro. Esta competencia entre sustitución y
eliminación recorre toda la síntesis orgánica. Este capítulo describe los
dos mecanismos de eliminación, su notable exigencia geométrica, la regla
que predice qué alqueno se forma y cómo orientar una reacción hacia uno u
otro lado.

\begin{recall}
La SN1 y la SN2 (\cref{ch:b1:nucleophilic-substitution}); las
\omterm{def:b1:stereochemistry-in-depth:representations}{proyecciones de Newman}, las \omterm{def:b1:stereochemistry-in-depth:isomers}{conformaciones} del ciclohexano y los
descriptores $Z$/$E$ (\cref{ch:b1:stereochemistry-in-depth}); los
\omterm{def:b1:electronic-effects:intermediates}{carbocationes} y las bases
(\cref{ch:b1:electronic-effects}).
\end{recall}

\section{El mecanismo E2}

\begin{definition}[\texorpdfstring{$\beta$}{Beta}-eliminación, E2 y E1]\label{def:b1:elimination:elimination}
Una \emph{$\beta$-eliminación}\index{beta-eliminacion@$\beta$-eliminación} elimina un \omterm{def:b1:electronic-effects:nucleophile}{grupo
saliente} Y de un carbono (el carbono $\alpha$) y un protón de un carbono
vecino (un carbono $\beta$), formando un enlace doble entre ellos. En el
\emph{mecanismo E2}\index{mecanismo E2} una base toma el protón mientras
se va Y, en una sola \omterm{def:b1:elementary-steps:elementary-step}{etapa elemental}; en el \emph{mecanismo
E1}\index{mecanismo E1}, primero se va Y, lo que da un \omterm{def:b1:electronic-effects:intermediates}{carbocatión}, que
después pierde un protón $\beta$.
\end{definition}

\begin{proposition}[Ley de velocidad de la E2]\label{prop:b1:elimination:e2-law}
Para la E2, $v = k[\mathrm{RY}][\mathrm{B}]$, donde B es la base.
\end{proposition}

\begin{proof}
Una sola \omterm{def:b1:elementary-steps:elementary-step}{etapa elemental} bimolecular, en la que intervienen el \omterm{def:b1:nucleophilic-substitution:substitution}{sustrato} y
la base: la ley de acción de masas de las \omterm{def:b1:elementary-steps:elementary-step}{etapas elementales}
(\cref{ch:b1:elementary-steps}).
\end{proof}

\begin{definition}[Antiperiplanar y sinperiplanar]\label{def:b1:elimination:periplanar}
Dos enlaces de átomos vecinos son \emph{antiperiplanares}\index{antiperiplanar}
cuando están en un mismo plano en lados opuestos (\omterm{def:b1:stereochemistry-in-depth:conformations}{ángulo de torsión} de
$180^\circ$), y \emph{sinperiplanares}\index{sinperiplanar} cuando están
en un mismo plano del mismo lado ($0^\circ$).
\end{definition}

\begin{omfigure}
\begin{tikzpicture}[font=\small, >={Stealth}]
  \pic (n) at (0,0) {omnewman={90}{270}};
  \node at (n-f1) {H}; \node at (n-f2) {\ce{CH3}}; \node at (n-f3) {H};
  \node at (n-b1) {Br}; \node at (n-b2) {\ce{CH3}}; \node at (n-b3) {H};
  \node (b) at (-2.6,2.0) {\ce{EtO-}};
  \draw[->, thick, omThm] (b) to[bend left=20] ($(n-f1)+(-0.15,0.1)$);
  \draw[->, thick, omThm] (0.12,0.8) to[bend left=40] (0.12,0.25);
  \draw[->, thick, omThm] (0.12,-0.6) to[bend right=35] ($(n-b1)+(0.25,0.15)$);
  \node[below] at (0,-2.0) {H y Br antiperiplanares};
  \draw[->, thick] (2.4,0) -- (3.6,0) node[midway, above] {E2};
  \node at (5.6,0) {\chemfig{C(-[:150]H_3C)(-[:210]H)=C(-[:30]H)(-[:-30]CH_3)}};
  \node at (5.6,-0.9) {\cip{E}-but-2-eno};
\end{tikzpicture}

{\small Reacción E2 del 2-bromobutano, vista a lo largo del enlace C3--C2
(C3 delante, C2, que lleva el Br, detrás) en una de sus dos \omterm{def:b1:stereochemistry-in-depth:isomers}{conformaciones}
\omterm{def:b1:elimination:periplanar}{antiperiplanares}; la otra, con el otro hidrógeno de C3 en anti respecto
del Br, da el \cip{Z}-but-2-eno. Tres pares de electrones se mueven a la
vez: la base toma el protón, el par C--H se convierte en el enlace $\pi$ y
el par C--Br se va con el bromo. Los grupos que se enfrentan en la
\omterm{def:b1:stereochemistry-in-depth:representations}{proyección de Newman} terminan del mismo lado del enlace doble.}
\end{omfigure}

\begin{proposition}[La E2 es anti y estereoespecífica]\label{prop:b1:elimination:anti}
La E2 transcurre desde la \omterm{def:b1:stereochemistry-in-depth:isomers}{conformación} en la que los enlaces C--H y C--Y
son \omterm{def:b1:elimination:periplanar}{antiperiplanares}. Es estereoespecífica: \omterm{def:b1:nucleophilic-substitution:substitution}{sustratos} \omterm{def:b1:stereochemistry-in-depth:enantiomers}{diastereoisómeros}
dan \omterm{def:b1:stereochemistry-in-depth:isomers}{estereoisómeros} distintos del alqueno, fijados por la disposición
anti.
\end{proposition}

\begin{proof}
En el \omterm{def:b1:elementary-steps:energy-profile}{estado de transición}, el \omterm{def:b1:lewis-resonance-vsepr:lewis}{par enlazante} C--H debe entrar en la
región antienlazante del enlace C--Y mientras los dos carbonos se vuelven
trigonales; el solapamiento es máximo cuando los dos enlaces son
paralelos, en un mismo plano. La disposición anti es alternada (de baja
energía) y mantiene la base lejos del \omterm{def:b1:electronic-effects:nucleophile}{grupo saliente}; la sin es
eclipsada. Con H e Y en anti, los otros dos grupos de cada carbono quedan
fijados: los que están del mismo lado en la \omterm{def:b1:stereochemistry-in-depth:representations}{proyección de Newman} terminan
en \textit{cis} en el enlace doble. Un diastereoisómero tiene dos grupos
intercambiados en un carbono, y da el otro alqueno.
\end{proof}

\begin{proposition}[E2 en un ciclohexano]\label{prop:b1:elimination:cyclohexane-e2}
En un anillo de ciclohexano, la E2 exige que el \omterm{def:b1:electronic-effects:nucleophile}{grupo saliente} y el
hidrógeno $\beta$ sean ambos \omterm{def:b1:stereochemistry-in-depth:conformations}{axiales} (\textit{trans}-diaxiales); un \omterm{def:b1:electronic-effects:nucleophile}{grupo
saliente} retenido en posición \omterm{def:b1:stereochemistry-in-depth:conformations}{ecuatorial} no puede reaccionar hasta que el
anillo se invierte.
\end{proposition}

\begin{proof}
En carbonos contiguos de una \omterm{def:b1:stereochemistry-in-depth:conformations}{silla}, dos enlaces \omterm{def:b1:stereochemistry-in-depth:conformations}{axiales} son
\omterm{def:b1:elimination:periplanar}{antiperiplanares} (uno hacia arriba y otro hacia abajo, paralelos al eje);
un enlace \omterm{def:b1:stereochemistry-in-depth:conformations}{ecuatorial} solo está en anti respecto de enlaces C--C del
anillo, nunca de un enlace C--H.
\end{proof}

\begin{omfigure}
\begin{tikzpicture}[font=\small]
  \pic (a) at (0,0) {omchair};
  \draw[thick, omThm] (a-c1) -- (a-a1) node[above] {Cl};
  \draw[thick] (a-c1) -- (a-e1) node[right] {H};
  \draw[thick, omDef] (a-c2) -- (a-a2) node[below] {H};
  \draw[thick, omDef] (a-c6) -- (a-a6) node[below] {H};
  \node[align=left, anchor=west] at (3.0,0.2) {Cl axial (arriba) en C1;\\H axial (abajo) en los dos vecinos:\\dos H antiperiplanares};
\end{tikzpicture}

{\small Disposición \textit{trans}-diaxial en una \omterm{def:b1:stereochemistry-in-depth:conformations}{silla}: el cloro \omterm{def:b1:stereochemistry-in-depth:conformations}{axial} es
\omterm{def:b1:elimination:periplanar}{antiperiplanar} respecto de los hidrógenos \omterm{def:b1:stereochemistry-in-depth:conformations}{axiales} de los dos carbonos
vecinos (en azul).}
\end{omfigure}

\section{El mecanismo E1}

\begin{proposition}[Ley de velocidad de la E1]\label{prop:b1:elimination:e1-law}
Para la E1, $v = k[\mathrm{RY}]$, pues la velocidad es la de la
ionización.
\end{proposition}

\begin{proof}
Como en la SN1 (\cref{prop:b1:nucleophilic-substitution:sn1-law}): el
\omterm{def:b1:electronic-effects:intermediates}{carbocatión} se forma en una etapa lenta y se consume deprisa, aquí por
pérdida de un protón hacia el disolvente o una \omterm{def:b1:predominant-reaction:strong-acid}{base débil}; la
\omterm{def:b1:elementary-steps:qssa}{aproximación del estado estacionario} deja $v = k_1[\mathrm{RY}]$.
\end{proof}

La E1 comparte su \omterm{def:b1:electronic-effects:intermediates}{carbocatión} con la SN1: los \omterm{def:b1:nucleophilic-substitution:substitution}{sustratos} terciarios en
\omterm{def:b1:intermolecular-forces-solvents:solvent-classes}{disolventes próticos} dan ambas, y el calentamiento favorece el alqueno.
Como el \omterm{def:b1:electronic-effects:intermediates}{carbocatión} es plano y el protón se pierde después, la E1 no
tiene exigencia geométrica y no es estereoespecífica; da sobre todo el
alqueno más estable, normalmente el \cip{E}.

\begin{omfigure}
\setchemfig{atom sep=2.1em}
\schemestart
\chemfig{C(-[2]CH_3)(-[4]H_3C)(-[6]CH_3)-Br}
\arrow{->[lenta]}
\chemfig{C^{+}(-[2]CH_3)(-[:210]H_3C)(-[:-30]CH_3)}
\arrow{->[rápida][\ce{-H+}]}
\chemfig{C(-[2]CH_3)(-[:210]H_3C)=[:-30]CH_2}
\schemestop
\par\vspace{4pt}

{\small E1: ionización del 2-bromo-2-metilpropano al \omterm{def:b1:electronic-effects:intermediates}{carbocatión} y,
después, pérdida de un protón de un grupo metilo hacia el disolvente, que
da el 2-metilpropeno.}
\end{omfigure}

\section{Regioselectividad: la regla de Zaitsev}

\begin{definition}[Reacciones regioselectivas y estereoselectivas]\label{def:b1:elimination:selectivity}
Una reacción que puede dar varios \omterm{def:b1:stereochemistry-in-depth:isomers}{isómeros constitucionales} pero forma
mayoritariamente uno es una \emph{reacción regioselectiva}\index{reaccion regioselectiva@reacción regioselectiva};
una que puede dar varios \omterm{def:b1:stereochemistry-in-depth:isomers}{estereoisómeros} pero forma mayoritariamente uno
es una \emph{reacción estereoselectiva}\index{reaccion estereoselectiva@reacción estereoselectiva}.
\end{definition}

\begin{proposition}[Regla de Zaitsev]\label{prop:b1:elimination:zaitsev}
Cuando pueden arrancarse varios hidrógenos $\beta$, el alqueno
mayoritario suele ser el más sustituido (el más estable), y el isómero
\cip{E} antes que el \cip{Z} (\emph{regla de Zaitsev}\index{regla de Zaitsev}).
Las bases voluminosas y las restricciones geométricas, como la exigencia
trans-diaxial, pueden imponerse a ella.
\end{proposition}

\begin{proof}
Los grupos alquilo sobre los carbonos de un enlace doble lo estabilizan
(como estabilizan los \omterm{def:b1:electronic-effects:intermediates}{carbocationes}), y los alquenos \cip{E} evitan el
abarrotamiento de los grupos \textit{cis}. Los \omterm{def:b1:elementary-steps:energy-profile}{estados de transición} de
la E2 y la E1 tienen carácter parcial de enlace doble, de modo que el
alqueno más estable se forma más deprisa. Una base voluminosa alcanza con
más facilidad los hidrógenos menos impedidos; y un hidrógeno que no puede
ser \omterm{def:b1:elimination:periplanar}{antiperiplanar} respecto del \omterm{def:b1:electronic-effects:nucleophile}{grupo saliente} no se elimina por E2 en
absoluto, sea cual sea la estabilidad del alqueno que daría.
\end{proof}

\begin{method}[Predecir el alqueno]\label{met:b1:elimination:predict-alkene}
\begin{enumerate}
  \item Se enumeran los carbonos $\beta$ que llevan hidrógenos.
  \item Para la E2, se conservan solo los hidrógenos que pueden ser
        \omterm{def:b1:elimination:periplanar}{antiperiplanares} respecto de Y (en un anillo: trans-diaxiales, en
        una \omterm{def:b1:stereochemistry-in-depth:conformations}{silla} que realmente pueda formarse).
  \item Entre ellos, se favorece el que da el alqueno más sustituido
        (Zaitsev), salvo si la base es voluminosa.
  \item Para cada uno, se dibuja la \omterm{def:b1:stereochemistry-in-depth:representations}{proyección de Newman} en la
        \omterm{def:b1:stereochemistry-in-depth:conformations}{conformación anti} para leer la geometría \cip{E}/\cip{Z}.
\end{enumerate}
\end{method}

\section{¿Sustitución o eliminación?}

\begin{proposition}[Sustitución frente a eliminación]\label{prop:b1:elimination:sn-vs-e}
La eliminación se ve favorecida por las \omterm{def:b1:predominant-reaction:strong-acid}{bases fuertes} y voluminosas, por
la temperatura alta y por los \omterm{def:b1:nucleophilic-substitution:substitution}{sustratos} terciarios o impedidos; la
sustitución, por los buenos \omterm{def:b1:electronic-effects:nucleophile}{nucleófilos} que son \omterm{def:b1:predominant-reaction:strong-acid}{bases débiles}, por la
temperatura baja y por los \omterm{def:b1:nucleophilic-substitution:substitution}{sustratos} primarios.
\end{proposition}

\begin{proof}
La fuerza de la base favorece el ataque al hidrógeno; el volumen
dificulta el ataque a un carbono impedido, pero no a un hidrógeno
expuesto. Un \omterm{def:b1:electronic-effects:carbon-class}{carbono terciario} es inaccesible a la SN2, y su \omterm{def:b1:electronic-effects:intermediates}{carbocatión}
pierde protones con facilidad. La eliminación convierte una molécula en
dos o tres (alqueno, \omterm{def:b1:electronic-effects:nucleophile}{grupo saliente}, base protonada): se ve favorecida al
subir la temperatura, un resultado que se justifica en el volumen del
segundo año.
\end{proof}

\begin{omfigure}
\begin{tabular}{llll}
\hline
\omterm{def:b1:nucleophilic-substitution:substitution}{sustrato} & \omterm{def:b1:predominant-reaction:strong-acid}{base fuerte}, pequeña & \omterm{def:b1:predominant-reaction:strong-acid}{base fuerte}, voluminosa & \omterm{def:b1:predominant-reaction:strong-acid}{base débil}, buen \omterm{def:b1:electronic-effects:nucleophile}{nucleófilo} \\
\hline
primario & SN2 (algo de E2) & E2 & SN2 \\
secundario & E2 (y SN2) & E2 & SN2 (aprótico), SN1/E1 (prótico) \\
terciario & E2 & E2 & SN1/E1 (prótico) \\
\hline
\end{tabular}

\smallskip
{\small Una primera orientación: la reacción principal en cada caso.
Calentar favorece la eliminación en cada columna.}
\end{omfigure}

\section{Ejercicios}

\begin{exercise}[$\star$]\label{exo:b1:elimination:1}
Dense los alquenos que pueden formarse por eliminación de HBr a partir
del 2-bromobutano, del 2-bromo-2-metilbutano y del bromociclohexano, y el
que la \omterm{prop:b1:elimination:zaitsev}{regla de Zaitsev} predice como mayoritario.
\end{exercise}

\begin{exercise}[$\star$]\label{exo:b1:elimination:2}
Escríbanse las \omterm{def:b1:rate-laws:order}{leyes de velocidad} de la E1 y de la E2. ¿Cómo se
distinguirían experimentalmente para el 2-bromo-2-metilpropano en etanol
con etóxido?
\end{exercise}

\begin{exercise}[$\star$]\label{exo:b1:elimination:3}
Dibújense las \omterm{def:b1:stereochemistry-in-depth:representations}{proyecciones de Newman} del 2-bromobutano a lo largo de
C2--C3 en sus tres \omterm{def:b1:stereochemistry-in-depth:conformations}{conformaciones alternadas}, y dígase cuáles pueden
sufrir una E2.
\end{exercise}

\begin{exercise}[$\star$]\label{exo:b1:elimination:4}
Predígase la reacción principal (SN1, SN2, E1, E2): el bromoetano con
hidróxido de sodio en agua a \qty{25}{\celsius}; el
2-bromo-2-metilpropano con terc-butóxido de potasio; el 2-bromopropano
con yoduro de sodio en propanona; el 2-bromo-2-metilpropano en etanol
caliente.
\end{exercise}

\begin{exercise}[$\star\star$]\label{exo:b1:elimination:5}
Demuéstrese, con \omterm{def:b1:stereochemistry-in-depth:representations}{proyecciones de Newman}, que la eliminación E2 de HBr a
partir de los dos \omterm{def:b1:stereochemistry-in-depth:enantiomers}{diastereoisómeros} del 2-bromo-3-metilpentano que pueden
perder el hidrógeno de C3 da los dos \omterm{def:b1:stereochemistry-in-depth:isomers}{estereoisómeros} del
3-metilpent-2-eno.
\end{exercise}

\begin{exercise}[$\star\star$]\label{exo:b1:elimination:6}
Escríbase el mecanismo de la deshidratación del 2-metilbutan-2-ol en
ácido sulfúrico concentrado caliente (E1) y dese el alqueno mayoritario.
\end{exercise}

\begin{exercise}[$\star\star$]\label{exo:b1:elimination:7}
Explíquese por qué el \textit{trans}-1-bromo-4-terc-butilciclohexano
reacciona muy despacio por E2, mientras que su isómero \textit{cis}
reacciona deprisa. (El grupo terc-butilo permanece \omterm{def:b1:stereochemistry-in-depth:conformations}{ecuatorial}.)
\end{exercise}

\begin{exercise}[$\star\star$]\label{exo:b1:elimination:8}
El 2-bromo-2-metilbutano con etóxido de sodio da sobre todo
2-metilbut-2-eno; con terc-butóxido de potasio, sobre todo
2-metilbut-1-eno. Explíquese.
\end{exercise}

\begin{exercise}[$\star\star$]\label{exo:b1:elimination:9}
¿Por qué nunca se observa eliminación con el bromometano? ¿Y con el
1-bromo-2,2-dimetilpropano por E2?
\end{exercise}

\begin{exercise}[$\star\star\star$]\label{exo:b1:elimination:10}
En etanol a \qty{25}{\celsius}, el 2-bromo-2-metilpropano da una mezcla
del éter (SN1) y del alqueno (E1) cuyas proporciones no dependen de la
concentración del \omterm{def:b1:nucleophilic-substitution:substitution}{sustrato}. Demuéstrese que ambos productos proceden del
mismo \omterm{def:b1:electronic-effects:intermediates}{carbocatión} y que su razón es la razón de dos \omterm{def:b1:rate-laws:order}{constantes de
velocidad}.
\end{exercise}

\begin{exercise}[$\star\star\star$]\label{exo:b1:elimination:11}
Uno de los \omterm{def:b1:stereochemistry-in-depth:enantiomers}{diastereoisómeros} del 1,2-dibromo-1,2-difeniletano (el meso)
da, por E2 con etóxido, un solo estereoisómero del
1-bromo-1,2-difenileteno. Hállese cuál, con una \omterm{def:b1:stereochemistry-in-depth:representations}{proyección de Newman}.
\end{exercise}

\begin{exercise}[$\star\star\star$]\label{exo:b1:elimination:12}
Con las energías del \cref{ch:b1:stereochemistry-in-depth}, explíquese
por qué una reacción E2 sobre un ciclohexano puede frenarse en un factor
de cientos cuando el \omterm{def:b1:electronic-effects:nucleophile}{grupo saliente} debe volverse \omterm{def:b1:stereochemistry-in-depth:conformations}{axial} en una \omterm{def:b1:stereochemistry-in-depth:conformations}{silla}
abarrotada. Exprésese la velocidad como el producto de una fracción de
confórmero y una \omterm{def:b1:rate-laws:order}{constante de velocidad}.
\end{exercise}

\section{Problema: los cloruros de mentilo y de neomentilo}

\begin{problem}[{Problema de fin de semana --- las sillas de dos cloruros
diastereoisómeros, los hidrógenos antiperiplanares disponibles en cada
uno, los alquenos formados y la razón de sus velocidades de E2}]
\label{pb:b1:elimination:1}
El cloruro de mentilo y el cloruro de neomentilo son los cloruros del
mentol y del neomentol (\cref{ch:b1:stereochemistry-in-depth}): en el
anillo, C1 lleva el Cl; C2, el grupo isopropilo, y C5, el grupo metilo. En
el cloruro de mentilo, los tres grupos pueden ser todos \omterm{def:b1:stereochemistry-in-depth:conformations}{ecuatoriales}; en
el cloruro de neomentilo, el Cl es \omterm{def:b1:stereochemistry-in-depth:conformations}{axial} cuando los dos grupos alquilo son
\omterm{def:b1:stereochemistry-in-depth:conformations}{ecuatoriales}. Ambos se tratan con etóxido de sodio en etanol, una reacción
E2. Datos del problema: el \qty{95}{\percent} del cloruro de neomentilo
está en la \omterm{def:b1:stereochemistry-in-depth:conformations}{silla} con el Cl \omterm{def:b1:stereochemistry-in-depth:conformations}{axial}; en el cloruro de mentilo, la \omterm{def:b1:stereochemistry-in-depth:conformations}{silla} con
el Cl \omterm{def:b1:stereochemistry-in-depth:conformations}{axial} (los tres grupos \omterm{def:b1:stereochemistry-in-depth:conformations}{axiales}) es una fracción \num{5.0e-3};
supóngase que cada hidrógeno \omterm{def:b1:elimination:periplanar}{antiperiplanar} de una \omterm{def:b1:stereochemistry-in-depth:conformations}{silla} reactiva
reacciona con la misma \omterm{def:b1:rate-laws:order}{constante de velocidad}. El cloruro de neomentilo da
un \qty{78}{\percent} del alqueno trisustituido.

\medskip
\noindent\textbf{Parte I --- Las \omterm{def:b1:stereochemistry-in-depth:conformations}{sillas}.}
\begin{enumerate}
  \item Dibújese la \omterm{def:b1:stereochemistry-in-depth:conformations}{silla} preferida del cloruro de mentilo.
  \item Dibújese la \omterm{def:b1:stereochemistry-in-depth:conformations}{silla} preferida del cloruro de neomentilo.
  \item ¿Son los dos cloruros \omterm{def:b1:stereochemistry-in-depth:enantiomers}{enantiómeros} o \omterm{def:b1:stereochemistry-in-depth:enantiomers}{diastereoisómeros}?
  \item Enúnciese la exigencia geométrica de la E2 en un anillo.
  \item ¿En qué \omterm{def:b1:stereochemistry-in-depth:conformations}{silla} puede reaccionar el cloruro de mentilo? Dibújese.
  \item ¿Por qué es rara esa \omterm{def:b1:stereochemistry-in-depth:conformations}{silla}? Úsese el coste energético de los
        grupos \omterm{def:b1:stereochemistry-in-depth:conformations}{axiales}.
\end{enumerate}

\medskip
\noindent\textbf{Parte II --- Los hidrógenos anti.}
\begin{enumerate}[resume]
  \item Enumérense los carbonos $\beta$ del cloruro (C2 y C6) y sus
        hidrógenos.
  \item En la \omterm{def:b1:stereochemistry-in-depth:conformations}{silla} reactiva del cloruro de neomentilo, ¿qué hidrógenos son
        \omterm{def:b1:elimination:periplanar}{antiperiplanares} respecto del Cl?
  \item En la \omterm{def:b1:stereochemistry-in-depth:conformations}{silla} reactiva del cloruro de mentilo, ¿es \omterm{def:b1:stereochemistry-in-depth:conformations}{axial} el hidrógeno
        de C2? ¿Es \omterm{def:b1:stereochemistry-in-depth:conformations}{axial} un hidrógeno de C6?
  \item ¿Cuántos hidrógenos \omterm{def:b1:elimination:periplanar}{antiperiplanares} ofrece cada cloruro?
  \item Dibújese una \omterm{def:b1:stereochemistry-in-depth:representations}{proyección de Newman} a lo largo de C1--C6 para la
        \omterm{def:b1:stereochemistry-in-depth:conformations}{silla} reactiva del cloruro de mentilo.
  \item ¿Por qué una eliminación \omterm{def:b1:elimination:periplanar}{sinperiplanar} no podría salvar la
        reacción?
\end{enumerate}

\medskip
\noindent\textbf{Parte III --- Los productos.}
\begin{enumerate}[resume]
  \item Nómbrense el alqueno que se forma al arrancar el hidrógeno de C2 y
        el que se forma al arrancar un hidrógeno de C6.
  \item ¿Cuál favorece la \omterm{prop:b1:elimination:zaitsev}{regla de Zaitsev}, y por qué?
  \item ¿Qué da el cloruro de neomentilo? ¿Es coherente la razón 78/22 con
        la \omterm{prop:b1:elimination:zaitsev}{regla de Zaitsev}?
  \item ¿Qué da el cloruro de mentilo? ¿Es coherente con la \omterm{prop:b1:elimination:zaitsev}{regla de
        Zaitsev}?
  \item ¿Qué reacción es estereoespecífica, regioselectiva o ambas cosas?
  \item ¿Por qué la sustitución (SN2) es minoritaria aquí con etóxido?
\end{enumerate}

\medskip
\noindent\textbf{Parte IV --- Las velocidades.}
\begin{enumerate}[resume]
  \item Escríbase la velocidad de cada reacción en función de la fracción
        de \omterm{def:b1:stereochemistry-in-depth:conformations}{silla} reactiva y del número de hidrógenos \omterm{def:b1:elimination:periplanar}{antiperiplanares}.
  \item Calcúlese la velocidad relativa del cloruro de neomentilo.
  \item Calcúlese la velocidad relativa del cloruro de mentilo.
  \item Calcúlese la razón de las velocidades
        $k(\text{neomentilo})/k(\text{mentilo})$.
\end{enumerate}
\end{problem}
